\section{Policy Gradient and PPO Methods}
\label{sec:policy_gradient_and_ppo_methods}

Policy-gradient methods optimize a parameterized policy directly by estimating the gradient of expected return with respect to the policy parameters.
For a stochastic policy $\pi_\theta(a\mid s)$, the objective can be written as $J(\theta)=\mathbb{E}_{\pi_\theta}[G_t]$, where $G_t$ denotes the discounted return from time $t$.
The policy-gradient estimator uses sampled trajectories to increase the log-probability of actions that produce positive advantage estimates:
\begin{equation}
    \nabla_\theta J(\theta)
    = \mathbb{E}_{\pi_\theta}\left[\nabla_\theta \log \pi_\theta(a_t\mid s_t) A^{\pi_\theta}(s_t,a_t)\right].
\end{equation}
In practice, $A^{\pi_\theta}$ is replaced by an empirical estimate computed from rewards and a learned value function.
This direct optimization is well suited to continuous action spaces because the policy can represent a probability distribution over real-valued actions.

Generalized Advantage Estimation (GAE) improves the bias--variance tradeoff of advantage computation by combining temporal-difference residuals over multiple horizons \parencite{schulman-2015}.
Let
\begin{equation}
    \delta_t = r_t + \gamma V_\varphi(s_{t+1}) - V_\varphi(s_t)
\end{equation}
be the one-step temporal-difference error for a value function $V_\varphi$.
GAE forms the advantage estimate
\begin{equation}
    \hat{A}_t = \sum_{l=0}^{\infty}(\gamma\lambda)^l \delta_{t+l},
\end{equation}
where $\lambda\in[0,1]$ controls the interpolation between low-variance one-step estimates and higher-variance Monte Carlo returns.
This mechanism is widely used in on-policy deep reinforcement learning because it stabilizes learning while allowing delayed rewards to influence earlier actions.

Proximal Policy Optimization (PPO) is a policy-gradient method designed to constrain each policy update so that optimization remains stable under minibatch stochastic gradient descent \parencite{schulman-2017}.
PPO uses the probability ratio
\begin{equation}
    \varrho_t(\theta)=\frac{\pi_\theta(a_t\mid s_t)}{\pi_{\theta_{\mathrm{old}}}(a_t\mid s_t)}
\end{equation}
between the new policy and the behavior policy that collected the data.
The clipped surrogate objective is
\begin{equation}
    L^{\mathrm{CLIP}}(\theta)=
    \mathbb{E}_t\left[
    \min\!\Bigl(\varrho_t(\theta)\hat{A}_t,
    \mathrm{clip}(\varrho_t(\theta),1-\epsilon,1+\epsilon)\hat{A}_t
    \Bigr)
    \right],
\end{equation}
where $\epsilon$ limits the incentive to move the new policy too far from the old policy.
The minimum prevents policy updates from obtaining large objective improvements by exceeding the clipping range in the direction favored by the advantage estimate.

PPO is commonly paired with a value-function loss and an entropy bonus.
The value loss trains the critic to predict returns, while the entropy term encourages stochasticity and discourages premature policy collapse.
The combined optimization problem is not a pure maximization of the clipped objective; it is a practical actor--critic training procedure that alternates between collecting on-policy rollouts and performing several optimization epochs over the collected batch.
Because the data become stale after the policy changes, PPO remains fundamentally on-policy and typically requires repeated environment interaction.

The use of PPO is relevant for intrinsic motivation because reward augmentation affects the advantage estimates used by the policy update.
If the reward is changed to
\begin{equation}
    r_t = r_t^{\mathrm{ext}} + \eta_t r_t^{\mathrm{int}},
\end{equation}
then the value targets and advantages reflect both task reward and intrinsic shaping.
A well-designed intrinsic signal can therefore guide the policy toward useful experience before extrinsic rewards are frequent, but a misleading signal can also dominate the optimization process.
PPO provides a stable backbone for comparing intrinsic rewards because the same policy optimizer can be used while varying only the exploration objective.